\part*{PPL2, Merkblatt zum Harvard-Konzept\\\normalsize
Severin Henzi, Simon Kronenberg, Malte Scheurer, Cyril Wendl}


Im Folgenden werden die Grundzüge des Harvard-Konzepts kurz präsentiert, danach werden Anwendungen mit Bezug zur Sekundarstufe II (Schweizer Bildungssystem) aufgezeigt.

\section*{Das Harvard-Konzept}
Im Buch ``Getting to Yes'' (Deutsch: ``Das Harvard-Konzept'') beginnen die Autoren \textcite{Fisher1981GettingTY} mit der Frage, welches die beste Art sei, wie Menschen mit ihren unterschiedlichen Arten umgehen können und der Feststellung, dass der Alltag von Familien, Paaren, Angestellten, Vorgesetzten, Konsumenten, Verkäufern und anderen Menschen, die in gemeinsamer Beziehung stehen, von der Herausforderung geprägt ist, in vielen Bereichen Einigungen erzielen zu können (``getting to yes''), ohne sich in Streitigkeiten zu verwickeln. Dabei wird häufig einer von zwei Ansätzen gewählt, um Einigungen zu erzielen: der ``harte'' Ansatz, bei dem jede Verhandlung zum Willens- und Überlebenskampf wird, oder der ``weiche'' Ansatz, bei dem eine Person nachgibt und Zugeständnisse macht, um einen Kompromiss zu erzielen. Die Autoren plädieren jedoch für einen anderen Ansatz, das sogenannte sachbezogene Verhandeln (``\textit{principled negotiation''}), welches sie im Rahmen des Harvard Negotiation Projects erarbeitet haben \cite{hnp}. Der grundlegende Befund dieses Ansatzes ist, dass Verhandlugen nicht als Präferenzen-basiertes ``Feilschen'' auszuführen seien, sondern vielmehr sachorientiert und basierend auf möglichst fairen Standards, welche zum grösstmöglichen gemeinsamen Nutzen führen. Aufgrund der Trennung vom Sachverhalt zu den Menschen lässt sich dieser Ansatz als ``hart'' in der Sache aber ``weich'' im Umgang mit Menschen bezeichnen. 

Die Autoren stellen drei Grundanforderungen an Verhandlungen: 
\begin{enumerate*}
\item Verhandlungen sollten zu einer Übereinkunft in einem gegebenen Thema führen;
\item Verhandlungen sollten effizient sein;
\item Verhandlungen sollten das Verhältnis zwischen beiden Parteien nicht stören, sondern wenn möglich stärken.
\end{enumerate*}

Darauf aufbauend entwickeln die Autoren vier Kriterien, um sachorientiertes Verhandeln zu definieren (s. \cref{table:gty})

\begin{table}[H]
\centering
\begin{tabular}{p{5cm}|p{5cm}||>{\columncolor[gray]{.9}}p{6cm}}
\multicolumn{2}{c||}{\textbf{Problemorientierte Ansätze}, ``Feilschen''} & \textbf{Lösungs-Orientierter Ansatz} \\
\textbf{Weich} & \textbf{Hart} & \textbf{Sachorientiert} \\\midrule\midrule
Teilnehmende sind Freunde & Teilnehmende sind Gegner & \textbf{1. Trennung von Person und Problem} \\\midrule
Zugeständnisse, um Beziehung zu erhalten & Forderungen durchsetzen & \textbf{2. Fokus auf Interessen, nicht Positionen} \\\midrule
Konflikte werden vermieden & Konflikte werden durchgesetzt & \textbf{3. Entwicklung von Optionen zum beiderseitigen Vorteil} \\\midrule
Auf Einigung Bestehen & Auf eigene Position bestehen & \textbf{4. Bestehen auf objektive Kriterien für Lösungen} \\
\end{tabular}
\caption{Prinzipien sachorientierten Handelns nach \cite[13]{Fisher1981GettingTY}}
\label{table:gty}
\end{table}

Punkt 1 bringt zum Ausdruck, dass das Problem auf bestmögliche Weise gelöst werden müsse, unabhängig persönlicher Interessen und Vertrauensverhältnisse. Punkt 2 besagt, dass Interessen aufgedeckt und diskutiert werden müssen, statt Positionen zu debattieren, und dass ein allzu schnelles Springen zur Schlussfolgerung vermieden werden sollte. Punkt 3 besagt, dass verschiedene Handlungs- oder Entscheidungsoptionen erarbeitet werden sollten, bevor direkt zur Entscheidung übergegangen wird. Punkt 4 fordert, dass ein Resultat basierend auf sachbezogenen Standards und möglichst unabhängig individueller Durchsetzungskraft erreicht werden. Statt also auf die Lösung oder auf Positionen zu fokussieren, soll immer wieder auf der Orientierung an objektiven Kriterien beharrte werden. Dem Druck anderer Teilnehmer, die auf eine Lösung drängen (weicher Ansatz) oder die eigene Position verteidigen (harter Ansatz) soll standgehalten werden.

\section*{Anwendungen im Kontext der Sekundarstufe II}
Im Kontext der Sekundarstufe II kommen vielfältige Anwendungsmöglichkeiten in den Sinn. Beispielsweise geschieht es fast täglich, dass einzelne Schüler von der Lehrperson gewisse Eingeständnisse einfordern: dabei kann es sich beispielsweise um die Wiederholung einer Prüfung handeln, um eine schlechte Note aufzubessern, um die Verschiebung von Abgaben im Falle von Krankheit oder um die Gewährung von Privilegien für allfällige ausserschulische Interessen. Zugleich verhandelt die Lehrperson auch regelmässig mit der Fachschaft, beispielsweise wenn es um die Anschaffung neuer Unterrichtsmaterialien oder um die Erstellung von Richtlinien geht, sowie auch mit dem Rektorat, wenn es um die eigene Karriere-Entwicklung oder Ressourcen-Allokation geht. Im Folgenden wird für einige dieser Beispiele analysiert, wie gewinnbringende Verhandlungsstrategien gemäss der Harvard-Theorie aussehen könnten, 

\begin{myExBox}[title=Anwendung 1: Verhandlung über Leistungsbeurteilung mit Schulrektorat]
(Fall Cyril)
\end{myExBox}


\begin{myExBox}[title=Anwendung 2: ...]

\end{myExBox}

\printbibliography